% Separate post-CP20 candidate. No promotion before full independent audits.
\subsection{The two possible anti-Cartan break patterns}

Retain the complete mixed-characteristic $(0,p)$, algebraically closed
residue, $p\ge5$, $e_K=p-1$, $\mu_p\subset K$ hypotheses of B58.
We consider the hypothetical different-character $p^2$ comparison
with irreducible residual inertia. By B57 and B58 only the full
wild image and residual projective order two remain. Write
\[
 W=Z\oplus U=M_2(\mathbf F_p),\quad q=p^2,\quad N=q-1,
 \quad A=q(p-1),\quad b_0=N(p-1).
\]
Here tame conjugation is $+1$ on the Cartan plane $Z$ and $-1$ on
the anti-Cartan plane $U$. Lemma~\ref{lem:B-RD58-final-Cartan}
identifies the last lower group as all of $Z$, with one break
$c>b_0$. Every anti-Cartan break is strictly smaller than $c$.

\begin{lemma}[The final-plane coefficient bound]
\label{lem:B-RD59-final-bound}
\status{Proved; full independent root and companion proof and exact-code audits passed.}
Fix one of the two models and a primitive residual coordinate
$z\ne0$. In its fixed uniformizer expansion, let $J$ be the first
nonzero degree of exact $p$-valuation one, and let $K$ be the first
nonzero degree prime to $p$, both below $qA$. An absent degree is
interpreted as $+\infty$. Then
\[
          J+pc\ge qA,\qquad K+c\ge qA,\qquad
                       \min(J+pc,K+c)=qA.              \tag{B-RD59.1}
\]
All degrees divisible by $q$ first change after $qA$ under a
nonzero element of $Z$.
\end{lemma}
\begin{proof}
The leading degree is $A$, with exact $p$-valuation two, and its
first change has degree $A+qc>qA$ because $c>b_0$.
Every higher degree divisible by $q$ changes later still.
The other two possible earliest changes have degrees $J+pc$ and
$K+c$. If their minimum were below $qA$, it could not occur
uniquely, since every nonzero Cartan matrix displaces every primitive
point by a nonzero first-level point. If both degrees attained that
minimum, their coefficient would be
\[
                           a d_g^p+b d_g,
\]
with at least one coefficient nonzero. The leading ramification
space $\{d_g:g\in Z\}$ has $q$ elements, because $Z$ is the final
lower group. A nonzero polynomial of degree at most $p$ cannot
vanish on this space. Thus both first changes are at least $qA$.
At least one must equal $qA$, as otherwise every change is later
than the required first-level displacement. This proves (B-RD59.1).
\end{proof}

\begin{lemma}[Both anti-Cartan patterns force an ordinary first term]
\label{lem:B-RD59-anti-patterns}
\status{Proved; full independent root and companion proof and exact-code audits passed.}
For each model the first nonzero prime-to-$p$ degree $K$ is
independent of $z\ne0$, has coefficient a fixed nonzero multiple
of $z$, and satisfies
\[
                           K+c=qA.                     \tag{B-RD59.2}
\]
Its change under every nonzero final Cartan element is the unique
first change. More precisely, if the anti-Cartan plane has a single
break $b$, then
\[
            0<b<b_0,\qquad K=A+Nb,\qquad
                           J\ge A+(q-p)b.              \tag{B-RD59.3}
\]
If it has two breaks $b_1<b_2$, then
\[
       J=A+(q-p)b_1,\qquad
       K=J+(p-1)b_2,\qquad b_1<b_0.                    \tag{B-RD59.4}
\]
These are conditional coefficient formulas, not existence assertions
for a formal group with either ramification pattern.
\end{lemma}
\begin{proof}
The lower filtration on the two-dimensional plane $U$ has either
one or two breaks. Its nonzero elements are all invertible as
endomorphisms of the residual space. Thus their displacement has
valuation $qA$ for every $z\ne0$.

Suppose first that all nonzero elements of $U$ have one break $b$.
Its leading coefficient space $D$ is an arbitrary two-dimensional
$\mathbf F_p$-subspace of $k$, not assumed to be a scalar multiple
of $\mathbf F_q$. Its root polynomial has the form
\[
                    P_D(X)=X^q+sX^p+tX,\qquad t\ne0.   \tag{B-RD59.5}
\]
Indeed, choose a basis $a,b$ of $D$ and put
$Y=X^p-a^{p-1}X$. Then
\[
 P_D(X)=Y^p-(b^p-a^{p-1}b)^{p-1}Y.
\]
The two factors giving its linear coefficient are nonzero by
linear independence. This proves (B-RD59.5) directly.
Tame conjugation on $U$ is $-1$, so $b\ne b_0$.
If $b>b_0$, the degree-$A$ change is already after $qA$.
Cancellation of an earlier change between the $J$ and $K$ terms
would require a polynomial of degree at most $p$ to vanish on
all of $D$, which is impossible. Thus $J+pb\ge qA$ and
$K+b\ge qA$, contradicting (B-RD59.1) because $c>b$.
Hence $b<b_0$.

As in the preceding argument, no change from $J$ or $K$ can occur
before $A+qb$: any such cancellation has degree at most $p$ in
the leading ramification coefficient. In degree $A+qb$, cancellation
must therefore be a nonzero scalar multiple of the full root
polynomial (B-RD59.5), with leading coefficient $(p-1)d_i z$.
Its nonzero linear coefficient forces $K+b=A+qb$, with
$Kf_{i,K}(z)$ a fixed nonzero multiple of $z$.
The possible $X^p$ term gives $J+pb\ge A+qb$; equality is
permitted and is not discarded. These are (B-RD59.3).
At the later Cartan break $c$ the prime-to-$p$ term is strictly
earlier than the other two classes, since
\[
 J+pc-(K+c)\ge(p-1)(c-b)>0,\qquad
 A+qc-(K+c)=N(c-b)>0.
\]
Equation (B-RD59.1) proves (B-RD59.2) and uniqueness.

Now suppose the two anti-Cartan breaks are $b_1<b_2<c$.
For the fixed coordinate $z$, consider the three possible first-change
degrees as functions of a lower break $b$:
\[
             A+qb,\qquad J+pb,\qquad K+b.               \tag{B-RD59.6}
\]
By (B-RD59.1), at least one of the last two degrees is strictly
less than $qA$ at $b_2$. Thus the minimum in (B-RD59.6) at
$b_2$ is below $qA$ and must occur at least twice.
It cannot include the first degree $A+qb_2$: if that degree
were at most both others at $b_2$, its strictly larger slope
$q$ would make it uniquely smallest at $b_1$, still below
$qA$. This would give an uncancelled earlier displacement.
Consequently
\[
             J+pb_2=K+b_2<A+qb_2,\qquad
                              J+pb_2<qA.               \tag{B-RD59.7}
\]
At $b_1$, the degree $J+pb_1$ is strictly smaller than $K+b_1$.
It is below $qA$. Cancellation now forces
$A+qb_1=J+pb_1$. These two equalities give (B-RD59.4), in
particular $b_1<b_0$.

Choose fixed nonzero anti-Cartan elements $g_1,g_2$ with these
two breaks. The cancellation at $b_1$ is between the leading
degree $A$ and degree $J$, so its coefficient equation makes
$f_{i,J}(z)$ a fixed nonzero multiple of $z$.
The cancellation at $b_2$ is then between degrees $J$ and $K$,
so it makes $f_{i,K}(z)$ a fixed nonzero multiple of
$f_{i,J}(z)$ and hence of $z$. The equations involve the same
fixed elements $g_1,g_2$ for every $z\ne0$. Thus the degrees
and scalar multiples do not depend on $z$.
Finally
\[
                  J+pc-(K+c)=(p-1)(c-b_2)>0.
\]
Together with (B-RD59.1), this gives (B-RD59.2); all degrees
divisible by $q$ are later by the first lemma. The prime-to-$p$
term is therefore uniquely first, proving the lemma in both cases.
\end{proof}

\begin{theorem}[Character recovery for every irreducible residual comparison]
\label{thm:B-RD59}
\status{Proved; full independent root and companion proof and exact-code audits passed.}
For $p\ge5$ and $e_K=p-1$, with $\mu_p\subset K$ and algebraically
closed residue field, isomorphic $p^2$-torsion determines the
requested special tame-order character for the potentially
supersingular height-two groups in the normalized comparison,
provided their common residual inertia representation is irreducible
over $\mathbf F_p$.
Consequently every remaining different-character pair in the general
minimal-ramification test has reducible residual inertia.
In particular this applies to any remaining failure at $25$-torsion
for $(p,e,h)=(5,4,2)$.
\end{theorem}
\begin{proof}
B57 and B58 leave only the projective-order-two full-wild case.
On its final Cartan plane the preceding lemma gives, for both models,
\[
                   d_i^q\beta_i(g)z=K f_{i,K}(z)d_g,
                     \qquad g\in Z.                   \tag{B-RD59.8}
\]
Since $f_{i,K}(z)$ is a fixed nonzero multiple of $z$, this
forces $\beta_i(g)=\kappa_i d_g$ on all of $Z$.
But $\beta_i:Z\to\mathbf F_q$ is bijective, while the
different-type residual intertwiner is Frobenius-semilinear and
therefore requires $\beta_2(g)=\beta_1(g)^p$.
As in B55 and B57, this would make Frobenius scalar-linear on
all of $\mathbf F_q$, a contradiction.
Thus no irreducible-residual different-character comparison remains.
\end{proof}

\begin{remark}[The unresolved residual case]
\label{rem:B-RD59-scope}
\status{Proved scope and limitation; reducible residual comparisons remain open.}
The theorem does not assume CM and does not replace a torsion module
by a different integral lattice. It uses the actual common division
field and the upper-ramification final-Cartan lemma of B58.
It leaves reducible residual representations entirely open, so the
unrestricted least-index interval at $(5,4,2)$ is still $\{1,2\}$.
The two-break argument compares all three slopes $q,p,1$; it is
not an extrapolation of the earlier cubic-wild two-slope argument.
The exact finite supplement
\path{code/reduction_descent/verify_tame_full_wild_quadratic.py}
checks 252 single-break and 124,788 two-break arithmetic cases,
27,648 three-slope candidates (432 satisfying the conditional ties),
and all 221 two-dimensional planes over $\mathbf F_5$, $\mathbf F_7$
and $\mathbf F_{11}$ in their cubic extensions. All these planes
have a nonzero middle coefficient in (B-RD59.5); the program checks
806 directly multiplied polynomial coefficients and 200,449 exact
root tests, as well as 364 scalar/Frobenius comparisons. The author
run passed in 1.1602564200002234 seconds. The complete independent
root and companion audits passed; their author-output-preserving
reruns took 1.100486028000887 and 1.1509096779991523 seconds,
respectively. The first version's
three-slope subtest had an empty admissible range; that source and
output are retained, and the revised range has the explicit
non-vacuity assertion just reported. These are finite checks of
the proof's identities, not constructions of a formal group or pair.
\end{remark}
